\begin{pastexam}[label={px:ntut-110-prog-03}]{北科大資工所 110 學年度 第 4 題}
（程式填空與追蹤）

\textbf{Problem 4 [11\%] [ 4-1 $\sim$ 4-3, each 2\%] [ 4-4 $\sim$ 4-8, each 1\%]}

Please trace the following C program and fill the blanks with correct statements or write down the
outputs of the statements.
\begin{minted}[fontsize=\small, breaklines]{c}
#include <stdio.h>
#include <stdlib.h>
typedef struct node { int data; struct node *next; } NODE;

int showTop(NODE* top) {
    return top->data;
}

void pushStack(_______ top, int data) {         /* Problem 4-1 (2 points) */
    NODE* node;

    node = (_______) malloc(sizeof(NODE));       /* Problem 4-2 (2 points) */
    if(node == NULL) {
        printf("Memory allocation failure!\n");
        exit(1);
    }
    node->data = data;
    node->next = *top;
    *top = node;
}

void showStack(NODE* top) {
    NODE* node = top;
    while(node != NULL) {
        printf("%d ", node->data);
        node = ___________;                      /* Problem 4-3 (2 points) */
    }
    printf("\n");
}

int popStack(_______ top) {                      /* Problem 4-4 (1 points) */
    int value;
    NODE* node = *top;
    if(node == NULL) {
        printf("Stack is empty!\n");
        exit(1);
    }
    value = (int) (node->data);
    _______ = node->next;                        /* Problem 4-5 (1 points) */
    free(node);
    return value;
}

int main(void) {
    NODE* stackTop=NULL;

    pushStack(&stackTop, 5);
    pushStack(&stackTop, 6);
    pushStack(&stackTop, 2);
    showStack(stackTop);                         /* Problem 4-6 (1 points) */
    printf("%d\n", popStack(&stackTop));         /* Problem 4-7 (1 points) */
    printf("%d\n", showTop(stackTop));           /* Problem 4-8 (1 points) */
    return 0;
}
\end{minted}
\end{pastexam}

\begin{pastsoln}
答案：
\begin{center}
\begin{tabular}{lll}
\toprule
4-1 & \texttt{NODE**} & 參數型別 \\
4-2 & \texttt{NODE*} & 轉型 \\
4-3 & \texttt{node->next} & \\
4-4 & \texttt{NODE**} & 參數型別 \\
4-5 & \texttt{*top} & \\
4-6 & \texttt{2 6 5} & 每個數後面一個空格，再換行 \\
4-7 & \texttt{2} & \\
4-8 & \texttt{6} & \\
\bottomrule
\end{tabular}
\end{center}

這是資結書第~6.1~節「用串列存」的堆疊（虛擬碼~6.2）：堆疊頂端就是串列的第一個節點，push 和 pop 都
在串列前端做，$O(1)$。

\textbf{4-1、4-4}：\texttt{main} 傳進來的是 \texttt{\&stackTop}，\texttt{stackTop} 的型別是
\texttt{NODE*}，取位址得到 \texttt{NODE**}。 函式裡也用 \texttt{*top} 讀寫 \texttt{main} 的
\texttt{stackTop}（\texttt{node->next = *top;}、\texttt{*top = node;}），所以參數必須是指向指標的指標。
如果寫成 \texttt{NODE* top}，函式只改到自己的副本，\texttt{main} 的 \texttt{stackTop} 永遠是
\texttt{NULL}。 pop 也要改頂端，同理填 \texttt{NODE**}。

\textbf{4-2}：\texttt{malloc} 回傳 \texttt{void*}，轉成要存的型別 \texttt{NODE*}，配置一個
\texttt{sizeof(NODE)} 大小的節點。

\textbf{4-3}：走訪串列，印完一個節點就往下一個走，\texttt{node = node->next;}，走到 \texttt{NULL}
停。

\textbf{4-5}：pop 時頂端改成原頂端的下一個節點，\texttt{*top = node->next;}，接著才 \texttt{free} 掉
原頂端。 順序不能反過來：先 \texttt{free(node)} 再讀 \texttt{node->next} 是讀已釋放的記憶體。

\textbf{4-6 到 4-8 的追蹤}（串列由頂端往下寫）：
\begin{center}
\begin{tabular}{ll}
\toprule
執行 & 堆疊（左邊是頂端） \\
\midrule
\texttt{pushStack(\&stackTop, 5)} & 5 \\
\texttt{pushStack(\&stackTop, 6)} & 6, 5 \\
\texttt{pushStack(\&stackTop, 2)} & 2, 6, 5 \\
\texttt{showStack} & 印出 \texttt{2 6 5}（每個數後接一個空格），換行 \\
\texttt{popStack} & 取出 2，剩 6, 5；印出 \texttt{2} \\
\texttt{showTop} & 頂端是 6，印出 \texttt{6} \\
\bottomrule
\end{tabular}
\end{center}
後進先出：最後 push 的 $2$ 最先印、最先被 pop。 \texttt{showStack} 的格式是 \texttt{"\%d "}，所以嚴格說
輸出是 ``\verb*|2 6 5 |'' 再換行，最後一個 $5$ 後面也有空格。
\end{pastsoln}
